\documentclass[10pt,a4paper]{article}
\usepackage[margin=18mm]{geometry}
\usepackage{amsmath,booktabs,array,tabularx}
\usepackage{cite,microtype}
\renewcommand{\citepunct}{], [}
\renewcommand{\citedash}{]--[}
\input{glyphtounicode}
\pdfgentounicode=1
\usepackage[hidelinks]{hyperref}
\hypersetup{pdftitle={An established traffic backbone for an elementary VSL task},
 pdfauthor={Differentiable traffic task review},
 pdfsubject={Literature comparison and methodological findings; first draft}}
\setlength{\parindent}{0pt}
\setlength{\parskip}{4pt}
\setlength{\tabcolsep}{4pt}
\renewcommand{\arraystretch}{1.14}
\newcommand{\head}[1]{\par\vspace{5pt}\noindent\textbf{#1}\par\nobreak\vspace{1pt}}
\newenvironment{points}{\begin{list}{$\bullet$}{%
 \setlength{\leftmargin}{1.3em}\setlength{\labelwidth}{0.6em}%
 \setlength{\labelsep}{0.4em}\setlength{\itemsep}{2pt}%
 \setlength{\parsep}{0pt}\setlength{\topsep}{2pt}}}{\end{list}}
\begin{document}
\begin{center}
{\large\bfseries An established traffic backbone for an elementary VSL task}\\[3pt]
Literature comparison, necessary adaptation and methodological findings\\[2pt]
First draft --- 21 October 2026
\end{center}
\vspace{-3pt}

\head{1. What a suitable backbone must provide}
Can the intended information and gradient questions be studied using an
established variable speed limit (VSL) model with very few changes? The
motivation is to keep
the traffic-model justification proportionate to an elementary verification
of a computational method. A backbone should supply the traffic state and
most evolution laws; simplicity concerns the assumptions and coupled rules
we must author and justify, rather than the number of references cited.

Consider a road with an ending lane E, its adjacent continuing lane A and a
farther continuing lane B. Two initial states can have identical E traffic
and identical combined A+B traffic at every road position, but different
A/B splits. We ask whether that concealed split changes merging service,
whether knowing it can improve a common speed decision, and whether complete
travel-time sensitivity can guide that decision. These are separate claims.
The objective is total system time, counting every initial unit until exit.
A suitable task must conserve traffic and respect shared receiving space.
Neither a plausible lane effect nor a differentiable formula establishes a
control advantage.

\head{2. What the closest published models offer}
The inspected literature offers two substantive routes: use an aggregate
VSL setting and ask about traffic location or history, or retain lane
densities to ask specifically about A/B loading. Table~1 compares what each
route already supplies and what our question would require.

\begin{center}
\textbf{Table 1. Published starting points and their fit to the question}\par
\vspace{3pt}
\begin{tabularx}{\linewidth}{@{}>{\raggedright\arraybackslash}p{0.19\linewidth}
 >{\raggedright\arraybackslash}p{0.37\linewidth}
 >{\raggedright\arraybackslash}X@{}}
\toprule
Published model & What is already supplied & Direct finding for this review\\
\midrule
Jin and Jin~\cite{jinjin}
& Spatial aggregate traffic and a link-average approximation; VSL entry
control, threshold capacity drop and upstream holding queues.
& The simplest inspected route if the hidden information is \emph{where}
traffic lies along the road. It has no A/B lane split; its discharge switch
still needs sensitivity qualification.\\[4pt]
Wei, Han and Wang~\cite{wei}
& Link traffic with entry/speed history, interacting speed cohorts and
capacity drop; advice acts on entering traffic.
& A serious history-aware aggregate alternative. Its native state omits
A/B loading, and its action timing differs from common advice throughout
our road.\\[4pt]
Nagalur Subraveti et al.~\cite{ns}
& Densities in successive lane cells, lane-change incentives and a lane-drop
application; congestion dynamics draw on Han et al.~\cite{han16}.
& Best fit among the inspected lane backbones for retaining the A/B question.
Add the VSL sending restriction~\cite{han17} and explicitly qualify movement
realization; direct source transcription is insufficient.\\[4pt]
Carey et al.~\cite{carey}
& Multilane cell transmission, intended lane changes, mandatory movements
and receiving-allocation alternatives.
& Useful support for conservative receiving and fallback rules. As the main
backbone, more of the VSL and permanent lane-end task would need specifying.\\
\bottomrule
\end{tabularx}
\end{center}

Jin and Jin therefore offer a genuine simplification, with a different
information contrast: equal total traffic can be distributed differently
along a road. A fair information study would use the same spatial traffic
model with different observations. Comparing their spatial and link-average
simulators would also change the dynamics. Likewise, lawful inference from
aggregate history must be admitted; aggregate models should not be reduced
to an arbitrary equal-lane guess.

For the chosen lane-loading question, the Subraveti framework supplies both
the missing state and reasons to redistribute it, including approaching the
lane end, density differences, keep-right behavior and cooperation. Its
Han lineage supplies a compatible starting point for congestion and VSL
budgets~\cite{han16,han17}. This is a bounded comparison of inspected methods,
not a proof that no simpler native model exists. The remaining issue is
whether the lane framework can be completed with a limited, explicit change.

\newpage
\head{3. How the selected model is completed}
\textbf{Concrete task and assumptions.}
The elementary member has ten 100\,m road sections, six-second updates and a
permanent three-to-two lane drop. A finite cohort is initialized inside the
road, independently of one constant common advice of 30--60\,km/h. There
are no later arrivals; crossing the common free A/B outlet line completes
a trip. E, A and B denote current lane stocks: early E-to-A traffic becomes
ordinary A traffic. The root assumptions are:
\begin{points}
\item Continuous mass, fixed geometry and a common lane flow--density law;
the time step respects sending and backward-wave storage bounds.
\item Local movements use pre-update traffic and space, advance one section,
and are applied simultaneously. New arrivals or freed space are not reused
within the same update.
\item Equal positive nominal priorities and unit receiving-space consumption
per unit of traffic; merging has no additional headway multiplier.
\end{points}

\textbf{Retained and adapted ingredients.}\par
Lane-density preferences retain the Subraveti framework~\cite{ns}; its
same-author thesis clarifies density lookahead~\cite{thesis}. Full Han
budgets bound available traffic and receiving service, with advice capping
sending~\cite{han16,han17}. Thus advice influences later merge service
through evolving traffic states. It does not directly prescribe a merge
fraction or change nominal priorities. Finite-domain, empty-lane and traffic-mode
conventions are specified explicitly for this member.

\textbf{The substantive replacement is movement realization.}
Review hand cases exposed donor/receiver budget conflicts in direct
transcription of the displayed native coupling; these are not claims about
recovered author code. The selected composition first normalizes all a lane's
requests against its one sending allowance. Requests into a destination
then share its one receiving allowance: small requests are satisfied, with
unused entitlement redistributed. Remaining traffic gets one opportunity
to continue in its own lane. At E's end there is no forward escape: it
enters A's outlet or waits. This adapts Carey receiving/fallback principles
and Wright priority sharing~\cite{carey,wright}. The complete composition is
ours; changing lane also advances traffic into the next road section.

Every realized transfer spends one donor budget and one receiver budget.
Subtracting and adding that same mass preserves conservation, while the
old-space bound protects storage. These joins make a complete model without
adding another physical mechanism. They do not imply maximum throughput,
equivalent congestion waves or inherited calibration. Calling this one
replacement component does not remove its assumptions. The closed-cohort,
one-advice comparison is also our task design, not an unchanged published
benchmark. The \href{vsl_joined_candidate_contract_and_qualification.pdf}{companion formulation}
retains the complete equations and hand cases.

\head{4. Recommendation and what the evidence establishes}
\textbf{Complete accounting and qualified sensitivity are established
analytically.} Under these assumptions, positive service and free outlets
eventually bring the cohort into an unconstrained regime. Remaining mass
need not become exactly zero in finite time: its distance to the exit gives
an exact remaining total system time. Consequently the objective counts
the full trip rather than discarding a small residual. For the fixed finite
model, the complete objective has smooth rational formulas between finitely
many possible advice switches. Its derivative includes earlier traffic
changes, residual-state changes and the direct residual travel-time effect.
At switches, jumps or unequal slopes may remain; derivatives can also be zero.
This is a fixed-step analytical result, not numerical gradient verification.

\textbf{Lane information changes service; its control benefit is unproved.}
The reviewed A/B pair has the same E and A+B profiles but different actual
merge service and exit flow. Advice is chosen before that informative
discharge; no identical hidden prehistory is asserted. The reduced observation loses predictive
information under this model. One extra A density in the last road section
identifies that pair completely, so it proves no need for all lane fields
or discrete order. Nor does different service establish different optimal
advice. In the sufficient small-cohort regime where all movements remain
unconstrained, equal traffic at each road position gives identical full
costs: maximum advice is optimal and lane information has zero action value.
The congested diagnostic lies outside that regime; its action gap is open.

\textbf{The resulting recommendation is conditional.}
Use the published lane framework, its congestion/VSL lineage and the explicit
movement replacement as a candidate for the selected information question.
Keep a model-aware E+(A+B) observer, the sufficient one-scalar augmentation,
maximum advice and equal-information scalar search as strong comparisons.
They separate information value from optimizer value. The thesis omits
bounded lane-change acceleration~\cite{thesis}; no earlier slow-entry/recovery
explanation or positive-action certificate transfers. Numerical correctness,
physical fidelity and useful-gradient performance remain future evidence.
The review establishes a coherent route with visible adaptation costs;
it does not establish that this route is globally simplest or best.

\typeout{REVIEW-MAIN-END-PAGE: \thepage}
\clearpage
\makeatletter
\renewcommand{\section}{\@ifstar{\head}{\head}}
\makeatother
\begin{thebibliography}{8}
\setlength{\itemsep}{4pt}\setlength{\parskip}{0pt}
\bibitem{jinjin} H.-Y. Jin and W.-L. Jin, ``Control of a lane-drop bottleneck
through variable speed limits,'' arXiv:1310.2658, 2013 (archived manuscript
dated Oct. 29, 2018). [Online]. Available:
\url{https://arxiv.org/abs/1310.2658}.
\bibitem{wei} L. Wei, Y. Han, and M. Wang, ``A first-order link-based flow
model with variable speed limits and capacity drops for freeway networks,''
\emph{Comput.-Aided Civ. Infrastruct. Eng.}, vol.~40, no.~21,
pp.~3200--3217, 2025, doi:
\href{https://doi.org/10.1111/mice.13492}{10.1111/mice.13492}.
\bibitem{ns} H.~H.~S. Nagalur Subraveti, V.~L. Knoop, and B. van Arem,
``First order multi-lane traffic flow model---an incentive based macroscopic
model to represent lane change dynamics,'' \emph{Transportmetrica B, Transp.
Dyn.}, vol.~7, no.~1, pp.~1758--1779, 2019 (archived article header: 2020),
doi: \href{https://doi.org/10.1080/21680566.2019.1700846}{10.1080/21680566.2019.1700846}.
\bibitem{han16} Y. Han, Y. Yuan, A. Hegyi, and S.~P. Hoogendoorn,
``New extended discrete first-order model to reproduce propagation of jam
waves,'' \emph{Transp. Res. Rec.}, vol.~2560, pp.~108--118, 2016, doi:
\href{https://doi.org/10.3141/2560-12}{10.3141/2560-12}.
\bibitem{han17} Y. Han, A. Hegyi, Y. Yuan, S. Hoogendoorn,
M. Papageorgiou, and C. Roncoli, ``Resolving freeway jam waves by discrete
first-order model-based predictive control of variable speed limits,''
\emph{Transp. Res. C, Emerg. Technol.}, vol.~77, pp.~405--420, 2017, doi:
\href{https://doi.org/10.1016/j.trc.2017.02.009}{10.1016/j.trc.2017.02.009}.
\bibitem{carey} M. Carey, C. Balijepalli, and D. Watling, ``Extending the
cell transmission model to multiple lanes and lane-changing,''
\emph{Netw. Spat. Econ.}, vol.~15, pp.~507--535, 2015, doi:
\href{https://doi.org/10.1007/s11067-013-9193-7}{10.1007/s11067-013-9193-7}.
\bibitem{thesis} H.~H.~S. Nagalur Subraveti, \emph{Lane-specific Traffic
Flow Control}, Ph.D. dissertation, Delft Univ. Technol., Delft, The
Netherlands, 2021, ch.~4, doi:
\href{https://doi.org/10.4233/uuid:e405a549-b1fe-4fa1-8090-4429612359d9}{10.4233/uuid:e405a549-b1fe-4fa1-8090-4429612359d9}.
\bibitem{wright} M.~A. Wright, G. Gomes, R. Horowitz, and
A.~A. Kurzhanskiy, ``On node models for high-dimensional road networks,''
\emph{Transp. Res. B, Methodol.}, vol.~105, pp.~212--234, 2017, doi:
\href{https://doi.org/10.1016/j.trb.2017.09.001}{10.1016/j.trb.2017.09.001}.
\end{thebibliography}
\end{document}
